% TOP_PROBLEM: {"problemId": "problem.brauers-kb-conjecture", "canonicalTitle": "Brauer's k(B) Conjecture", "releaseRank": 352, "primaryDomain": "algebra_representation_category", "primaryDomainLabel": "Algebra, representation & category theory", "displayStatus": "open_with_solved_subcases", "releaseStatus": "open_with_solved_subcases"}
% !TeX program = lualatex
% Definition and sources only; this file is not an LLM solution attempt.
\documentclass[11pt]{article}
\usepackage[margin=25mm]{geometry}
\usepackage{fontspec}
\setmainfont{Latin Modern Roman}
\usepackage{amsmath,amssymb,mathtools}
\usepackage{unicode-math}
\setmathfont{Latin Modern Math}
\usepackage{xurl,hyperref}
\hypersetup{colorlinks=true,urlcolor=blue}
\setcounter{secnumdepth}{0}
\setlength{\parindent}{0pt}
\setlength{\parskip}{5pt}
\setlength{\emergencystretch}{3em}
\begin{document}
\section*{\texorpdfstring{Brauer's k(B) Conjecture}{Brauer's k(B) Conjecture}}\label{352}
\subsection{Definitions and mathematical statement}
Fix a finite group $G$ and a prime $p$. Its ordinary irreducible complex characters split into $p$-blocks via modular representation theory over a splitting $p$-modular system. For a block $B$, write $\operatorname{Irr}(B)$ for its ordinary characters and $D$ for a defect group, defined up to conjugacy by the corresponding maximal Brauer pair. Brauer's conjecture asserts
\[
 k(B):=|\operatorname{Irr}(B)|\le |D|.
\]
The inequality applies to every finite group, every prime, and every block; $D$ need not be abelian. It counts ordinary characters, not the generally different number $\ell(B)$ of irreducible modular representations. No equality criterion or specified correspondence between characters and elements of $D$ is required.
\par\smallskip\textit{Formulation and concept references:} \href{https://arxiv.org/pdf/1911.10710}{[S1]}.

\subsection{Short English statement}
Can the number of ordinary irreducible characters in any block ever exceed the order of its defect group?
\subsection{Sources}
\begin{itemize}
\item[S1]Cesare Giulio Ardito and Benjamin Sambale, "Cartan matrices and Brauer's k(B)-Conjecture V," arXiv:1911.10710 (2019).
\url{https://arxiv.org/pdf/1911.10710}.
\textit{Formulation source.}
\end{itemize}
\end{document}
